% Source: https://artofproblemsolving.com/community/c3h3166500_2021_lmt_fall_guts_round_p25p36_lexington_mathematical_tournament
% Stable identity: aops-topic-3166500-q30
The final condition forces $d=e$. Write the larger of each successive ratio and its reciprocal as $u,v,w>1$, with $u+u^{-1}=7$, $v+v^{-1}=5$, $w+w^{-1}=3$. Each of $\dfrac{a}{b},\dfrac{b}{c},\dfrac{c}{d}$ independently equals its corresponding value or reciprocal. Inverting all three ratios does not change $f(e,a)$, leaving four candidate values. They are distinct: put $\alpha=\ln u$, $\beta=\ln v$, $\gamma=\ln w$, so $\alpha>\beta>\gamma>0$. The four values are $2\cosh(\alpha+\beta+\gamma)$ and the three values with one of these signs negative. Their absolute arguments are pairwise distinct; equality of two would force two of $\alpha,\beta,\gamma$ equal or one zero. Expanding $(u+u^{-1})(v+v^{-1})(w+w^{-1})$ groups its eight terms into exactly these four reciprocal-pair sums. Hence the requested sum is $7\cdot5\cdot3=\boxed{105}$.
